\begin{summarybox}
1. \textbf{四种命题结构}：
   \begin{itemize}
       \item 原命题：若 \(p\) 则 \(q\)；
       \item 逆命题：若 \(q\) 则 \(p\)；
       \item 否命题：若 \(\lnot p\) 则 \(\lnot q\)；
       \item 逆否命题：若 \(\lnot q\) 则 \(\lnot p\)。
   \end{itemize}
2. \textbf{等价关系}：原命题 \(\Leftrightarrow\) 逆否命题；逆命题 \(\Leftrightarrow\) 否命题。
3. \textbf{记忆口诀}：逆命题交换位置，否命题两边否定，逆否命题交换且否定。
4. \textbf{应用场景}：用于逻辑推理、间接证明（反证法实质是证明逆否命题）、判断命题真假。
5. \textbf{易错点}：区分"否命题"与"命题的否定"；注意否定词的正确写法。
\end{summarybox}